% GDZ PPN236018647,native562;HKUST563. Complete88–90.
\clearpage\SourceLeaf{562}\ArithmeticHead{533}{R}{1798 IUN.—1798 IUN. SIVE IUL.}
\NumberedItem{88}
\indent Calculus probabilitatis contra \textsc{La Place} defensus.
\par\noindent\hfill Gott[ingae, 1798] Iun. 17.

\medskip\begingroup\small
\indent Auf diese Aufzeichnung nimmt \textsc{Gauss} in zwei Briefstellen unmittelbar Bezug, nämlich in dem Briefe an \textsc{Olbers} vom 24. Januar 1812 (abgedruckt Werke VIII, S. 140) und in dem an \textsc{Laplace} vom 30. Januar 1812 (abgedruckt S. 371 dieses Bandes). An \textsc{Olbers} schreibt \textsc{Gauss}: »Unter meinen Papieren finde ich, daß ich im Junius 1798 .\ .\ .\ .\ .\ .\ . zuerst \textsc{Laplaces} Methode gesehen und die Unverträglichkeit derselben mit den Grundsätzen der Wahrscheinlichkeitsrechnung in einem kurzen Notizen-Journal über meine mathematischen Beschäftigungen angezeigt habe« und in dem Briefe an \textsc{Laplace} heißt es (siehe oben S. 373): »le mois de Juin 1798 est l’époque où je l’ai rapprochée [nämlich die Methode der kleinsten Quadrate] aux principes du calcul des probabilités«. Daraus geht hervor, daß es sich in unserer Aufzeichnung um die von \textsc{Gauss} im art. 186. der \textit{Theoria motus} (1809, Werke VII, 1906, S. 254) besprochene Methode von \textsc{Laplace-Boscovich} handelt, die \textsc{Laplace} im XI. Abschnitt seiner Abhandlung \textit{Sur quelques points du Système du monde}, Histoire de l’Académie des Sciences, Année 1789, Paris 1793, Mémoires etc. S. 32 auseinandersetzt, und daß \textsc{Gauss} diese Abhandlung im Juni 1798 kennen gelernt hat. Man vergl. auch die Einleitung des Aufsatzes von A. \textsc{Galle} »\textsc{Gauss} als Geodät«, Werke XI 2.
\par\noindent\hfill\textsc{Klein. Schlesinger.}
\par\endgroup

\NumberedItem{89}
\indent Problema eliminationis ita solutum, ut nihil amplius desiderari possit.
\par\noindent\hfill Gott[ingae, 1798] Iun.

\medskip\begingroup\small
\indent Mit der Eliminationstheorie, auf die sich auch schon die Tagebuchnotiz Nr. 36 vom 16. September 1796 bezieht, hat sich \textsc{Gauss} beschäftigt, als er die Beweise seiner Vorgänger \textsc{Euler} und \textsc{Lagrange} für den Fundamentalsatz der Algebra einer kritischen Prüfung unterzog. (Vgl. die Tagebuchnotiz Nr. 80 vom Oktober 1797.) Im art. 8. der \textit{Dissertation} (1799, Werke III, S. 13) lesen wir bei der Besprechung des \textsc{Euler}schen Beweises: »Ceterum operae pretium esse videtur, in formulas illas, quae \(\alpha,\beta\) etc. rationaliter per \(u,B,C\) etc. exprimant, profundius et generalissime inquirere; de qua re aliisque ad eliminationis theoriam (argumentum haudquaquam exhaustum) pertinentibus alia occasione fusius agere suscipiam«. Über \textsc{Lagrange} und seine Kritik des \textsc{Euler}schen Beweises spricht sich \textsc{Gauss} im art. 12. der \textit{Dissertation} (Werke III, S. 20) aus. Auf diesen Punkt bezieht sich auch die Bemerkung von \textsc{Pfaff} in dem oben abgedruckten Briefe vom 8. Juli 1799, S. 104 unten und S. 105 dieses Bandes.
\par\noindent\hfill\textsc{Loewy.}
\par\endgroup

\NumberedItem{90}
\indent Varia elegantiuscula circa attractionem sphaerae.
\par\noindent\hfill[1798 Iun. sive Iul.]
